\documentclass[
    language=english,
    doctype=blog-post,
    institution=none,
]{../../omnilatex}

\RequirePackage{config/document-settings}
\RequirePackage{listings}
\RequirePackage{xcolor}

\lstset{
    basicstyle=\ttfamily\small,
    keywordstyle=\color{blue}\bfseries,
    commentstyle=\color{gray},
    stringstyle=\color{red},
    breaklines=true,
    frame=single,
    language=Python,
    captionpos=b,
}

\begin{document}

\maketitle

\begin{abstract}
    Designing a REST API that can handle a million requests per day requires more
    than just choosing the right HTTP verbs. In this post, I share the architecture
    decisions, performance pitfalls, and practical lessons our team learned while
    scaling a Django-based API from prototype to production.
\end{abstract}

\section{Introduction}

When we started building the data ingestion API at CloudForge, we expected modest
traffic: a few hundred requests per minute from internal tools. Six months later,
external partners were hitting our endpoints at twenty times that volume, and our
response times had ballooned from 45ms to over two seconds.

This post documents the architectural changes that brought our p95 latency back
under 100ms---without a complete rewrite.

\section{The Starting Point}

Our initial stack was straightforward: Django REST Framework behind an Nginx
reverse proxy, with PostgreSQL as the sole datastore. The core endpoint looked
something like this:

\begin{lstlisting}[caption={Original endpoint handler}]
class IngestionView(APIView):
    def post(self, request):
        serializer = DataPayloadSerializer(data=request.data)
        serializer.is_valid(raise_exception=True)

        # Synchronous database write
        record = DataRecord.objects.create(
            payload=serializer.validated_data["payload"],
            source=serializer.validated_data["source"],
            timestamp=timezone.now(),
        )

        # Synchronous downstream processing
        process_record(record)

        return Response(
            {"id": record.id, "status": "accepted"},
            status=status.HTTP_202_ACCEPTED,
        )
\end{lstlisting}

This worked fine at low volumes. The problem was \texttt{process\_record}, which
performed a chain of synchronous operations: validation, transformation, index
updates, and cache invalidation.

\section{Identifying the Bottleneck}

We instrumented the endpoint with timing metrics and discovered that 80\% of
request time was spent inside \texttt{process\_record}. The chain averaged 320ms
per request, and with a single worker process, requests queued behind one another.

\begin{table}[h]
    \centering
    \begin{tabular}{lrr}
        \toprule
        \textbf{Operation} & \textbf{Avg Time (ms)} & \textbf{Share} \\
        \midrule
        Deserialization     & 12   & 3\%   \\
        Database write      & 35   & 10\%  \\
        Validation          & 28   & 8\%   \\
        Transformation      & 95   & 28\%  \\
        Index update        & 78   & 23\%  \\
        Cache invalidation  & 62   & 18\%  \\
        Network overhead    & 32   & 10\%  \\
        \bottomrule
    \end{tabular}
    \caption{Breakdown of request processing time before optimization.}
\end{table}

\section{The Solution: Event-Driven Offloading}

The key insight was that most of the work inside \texttt{process\_record} did not
need to happen synchronously. Clients needed confirmation that their payload was
\emph{received}, not that it was fully processed.

We introduced a message queue (Redis Streams) and restructured the handler:

\begin{lstlisting}[caption={Refactored endpoint with async processing}]
class IngestionView(APIView):
    def post(self, request):
        serializer = DataPayloadSerializer(data=request.data)
        serializer.is_valid(raise_exception=True)

        record = DataRecord.objects.create(
            payload=serializer.validated_data["payload"],
            source=serializer.validated_data["source"],
            timestamp=timezone.now(),
            status="pending",
        )

        # Publish to queue for async processing
        stream_publish("ingestion:incoming", {
            "record_id": str(record.id),
            "source": serializer.validated_data["source"],
            "received_at": timezone.now().isoformat(),
        })

        return Response(
            {"id": record.id, "status": "accepted"},
            status=status.HTTP_202_ACCEPTED,
        )
\end{lstlisting}

A separate worker service consumed from the queue and handled the heavy lifting
independently. This decoupled ingestion throughput from processing throughput.

\section{Results}

After deploying the changes, we observed immediate improvements:

\begin{itemize}
    \item Median response time dropped from 380ms to 18ms.
    \item Throughput increased from 150 requests/second to 4,200 requests/second.
    \item Error rates fell below 0.01\% as backpressure from slow downstream
          services no longer affected the ingestion path.
\end{itemize}

\section{Conclusion}

The biggest performance win came not from caching, database optimization, or
hardware upgrades---it came from rethinking what work truly needed to be
synchronous. If your API is bottlenecked on processing, ask whether clients
actually need the result immediately. In our case, they did not, and the async
offload cut latency by 95\%.

\section*{Tags}

\texttt{\#REST} \texttt{\#API} \texttt{\#SoftwareArchitecture} \texttt{\#Scalability}
\texttt{\#Django} \texttt{\#Microservices}

\section*{Comments}

\begin{quote}
    \textbf{Marcus T.} --- Great write-up! We faced a similar issue with our
    notification service and switched to a fan-out pattern. Did you consider
    Kafka instead of Redis Streams?
\end{quote}

\begin{quote}
    \textbf{Elena Vasquez (Author)} --- We evaluated Kafka but the operational
    overhead was not justified at our scale. Redis Streams fit our deployment
    footprint. If we cross 100k messages/second sustained, I would reconsider.
\end{quote}

\begin{quote}
    \textbf{Priya N.} --- The timing breakdown table is really helpful for
    making the case to management. Bookmarking this for our next sprint planning.
\end{quote}

\end{document}
